\documentclass[a4paper,fleqn]{cas-sc}
\usepackage[UTF8]{ctex}

\usepackage[authoryear,longnamesfirst]{natbib}
\usepackage{amsmath,amssymb,amsthm,mathtools}
\usepackage{booktabs,array,tabularx,longtable}
\usepackage{enumitem}
\usepackage{xcolor}
\usepackage{graphicx}
\usepackage{microtype}
\usepackage[section]{placeins}

% XeLaTeX compatibility for the Elsevier CAS/STIX class: restore ASCII math glyph slots.
\mathcode`\-="202D
\mathcode`\*="202A

\setlist[itemize]{leftmargin=2em,itemsep=0.25em,topsep=0.35em}
\setlist[enumerate]{leftmargin=2.3em,itemsep=0.3em,topsep=0.35em}
\emergencystretch=1em
\makeatletter
\setlength{\@fptop}{0pt}
\makeatother

\newtheorem{theorem}{定理}[section]
\newtheorem{proposition}[theorem]{命题}
\newtheorem{lemma}[theorem]{引理}
\newtheorem{corollary}[theorem]{推论}
\theoremstyle{definition}
\newtheorem{definition}[theorem]{定义}
\newtheorem{assumption}[theorem]{假设}
\theoremstyle{remark}
\newtheorem{remark}[theorem]{注}

\newcommand{\M}{\mathcal M}
\newcommand{\Pth}{\mathcal P}
\newcommand{\Led}{\mathsf{Led}}
\newcommand{\R}{\mathbb R}
\newcommand{\Z}{\mathbb Z}
\newcommand{\N}{\mathbb N}
\newcommand{\dd}{\,\mathrm d}
\newcommand{\TL}{T_L}
\newcommand{\TR}{T_R}
\newcommand{\AL}{A_L}
\newcommand{\AR}{A_R}
\newcommand{\Ker}{\operatorname{ker}}
\newcommand{\Hom}{\operatorname{Hom}}
\newcommand{\supp}{\operatorname{supp}}
\newcommand{\adeg}{\operatorname{adeg}}
\newcommand{\1}{\mathbf 1}
\newcommand{\cone}{\operatorname{cone}}
\newcommand{\spanr}{\operatorname{span}_{\R}}
\newcommand{\codim}{\operatorname{codim}}
\newcommand{\orcidurl}[1]{\href{https://orcid.org/#1}{#1}}
\ExplSyntaxOn
\RenewDocumentCommand \orcidauthor { m m }
{
  \seq_gput_right:Nn \g_stm_orcid_seq
  {
    { \ttfamily \orcidurl{#1} }
    \parsename{#2}
    \space(\eadauthor)
  }
}
\ExplSyntaxOff

\begin{document}
\shorttitle{恢复路径历史}
\shortauthors{Cheng}

\title[mode=title]{恢复路径历史：精确通行压缩与可加选择}

\author[1]{Yushang Cheng}[orcid={0009-0001-3218-6423}]
\cormark[1]
\ead{1302573502@qq.com}

\affiliation[1]{organization={西南交通大学希望学院},
                city={成都},
                country={中国}}

\cortext[1]{通讯作者}

\begin{abstract}
本文研究一维目标相对调整历史中，哪些特征能够在一种可加、对时间顺序进行粗粒化的评价中保留下来。通行账户记录所处一侧、运动方向、跨越的状态水平、重复次数以及目标达成。任何对弱时间变换不变且对拼接可加的赋值，都通过一个通用通行载体分解；该载体的精确纤维则可由实际历史之间的可行变换刻画。在一维情形下，有序端点与总跨越剖面共同即可重构完整的有向通行账户。

平衡通行循环一致性（BPCC）在并非对加法封闭的路径域上给出一个精确的偏可加次序。成对消失基准一致性（PVBC）刻画了何时实际历史之间的比较能够由一族归一化的实数支撑赋值作精确的一致同意表示，同时无需把形式差分强行全序化。穿孔揭示次序闭性（PROC）进一步给出正则性条件：在该条件下，支撑族可表示为四个有向 Stieltjes 测度，而目标达成仍可保留一个边界原子。于是，每个正则赋值都可分离为端点进展、总通行调整和目标达成三个部分。

这一套路径几何还支持可控的经济学变化。给定有限多个具有共同端点的可行计时遍历，只要共同时间跨度足够长，就能仅通过在端点加入停留来同时匹配持续时间以及某个能够区分两端点的连续流量统计量的累计值，而不必改变遍历速度，也不改变其通行账户。若两个方案对具有相同的通行账户差分，则序数排名可跨环境转移；有限个已观察的方案选择还能对未观察的通行对比给出锐界。唯一共同标量化属于另一个余维为一的识别特例。最后，即使菜单的备选项数量固定，完整可加一致性审计所需的重复系数仍可能无界。
\end{abstract}

\begin{keywords}
调整历史 \sep 通行压缩 \sep 可加测量 \sep 有向通行 \sep 目标达成 \sep 方案选择
\end{keywords}

\begingroup\hfuzz=120pt\maketitle\endgroup
\hypersetup{
  hypertexnames=false,
  pdftitle={恢复路径历史：精确通行压缩与可加选择},
  pdfauthor={Yushang Cheng},
  pdfsubject={目标相对调整历史上的精确通行压缩与可加选择},
  pdfkeywords={调整历史, 通行压缩, 可加测量, 有向通行, 目标达成, 方案选择}
}

\noindent\textbf{JEL 分类：} C60; D11; D90; I12.


% Main body is split for archival readability on GitHub.
\section{引言}

设想小明正在分阶段减烟方案之间进行选择。两个方案可以具有相同的初始状态、终止状态、持续时间和累计暴露指标，却在是否反向、跨越哪些状态水平以及是否到达治疗目标等方面不同。本文的问题是：对于一个可加的通行评价而言，这类历史中的哪些信息可以无损压缩？这种压缩又会对可观察选择施加什么限制？

一条历史首先生成一个\emph{通行账户}。该账户记录所处一侧、方向、跨越的状态水平、重复次数以及目标达成。它有意忽略日历时间，也忽略那些在不改变有向跨越记录的情况下可以交换的闭合游程之先后顺序。这种抽象不是直接假定的，而是被刻画出来的：任何对弱时间变换不变且对路径拼接可加的赋值，都必须通过一个通用载体分解；定理~\ref{thm:fiber} 则证明其逆向的纤维刻画。两条有限步历史具有相同账户，当且仅当它们可以通过可行的时间变换、停留、同基点闭合游程交换以及闭合路径的循环旋转相互连接。因此，该定理精确说明了账户究竟遗忘了哪些时间顺序信息。

一维几何还带来第二个压缩结果。净有向通行由有序端点决定，因此有序端点与总跨越剖面共同即可重构完整的有向账本，包括目标处的轨迹信息。这给出一个更小的充分记录，并说明：一旦同时观察到端点流量与总通行，方向就不再是独立的原始量。

通行账户本身不包含偏好。可加测量从实际历史之间的比较出发，而这个实际历史域并不对任意形式加法封闭。平衡通行循环一致性（BPCC）排除这样一类平衡的揭示改进列表：其中账户差分总和恰好抵消但至少有一项严格改进；由此得到一个精确的偏可加次序。成对消失基准一致性（PVBC）进一步排除相对于一个固定全支撑基准呈无穷小的实际严格比较。它使实际历史次序可以由一族归一化实数支撑赋值的一致同意来精确表示，同时仍允许未揭示的形式方向保持偏序。

正则性是对上述表示的进一步细化。穿孔区间 $(a,b_n]$ 在 $b_n\downarrow a$ 时从通行记录中消失，而区间 $[0,b_n]$ 始终包含一次目标事件。穿孔揭示次序闭性（PROC）要求揭示比较对第一类消失扰动保持闭合。于是，正则支撑赋值诱导出四个有向 Stieltjes 测度，并允许在目标处存在一个原子。面向目标的方向单调性则为四个模块提供应用层面的方向性解释。

对于每个正则支撑赋值 $\Lambda$，一维跨越守恒给出
\[
 U_\Lambda(\gamma)
 =P_\Lambda(\gamma(1))-P_\Lambda(\gamma(0))
  +\sum_s\int V_s^\gamma\,\dd\kappa_{s,\Lambda}
  +\sum_s\alpha_{s,\Lambda}A_s(\gamma).
\]
这一分解主要用于解释：端点进展被从对方向反转保持偶对称的总通行部分以及目标达成部分中分离出来。当多个支撑赋值都与原始次序相容时，这些分量依赖于具体赋值，因此可观察结论应以序数形式或集合识别形式表述。

通行账户的经济学用途并不要求所有背景后果在任意历史之间都字面不变。命题~\ref{prop:matched-design} 构造了一类匹配方案的时间安排，而无需允许任意快的遍历。对于任意有限个具有共同端点、且存在应用上可行计时遍历的右侧方案，只要共同时间跨度足够长，允许在端点停留就足以同时匹配持续时间以及一个在两端点取不同值的连续流量统计量的累计值，同时保持每条实际遍历及其通行账户不变。因此，可以通过实验设计固定端点、持续时间以及某个累计暴露量，而让通行几何发生变化。

由此得到两个可观察含义。第一，相同的通行账户差分可以把原始序数排名跨方案环境转移。第二，在有限经验网格上，有限个已观察选择定义了一个可加补全多面体，其线性像为任意未观察通行对比给出锐界。货币计价物可以把同一有限设计中的界校准为货币单位，但真正的识别对象是与数据相容的通行对比区间，而不是人为选定的一个标量效用。

数学工具主要来自可加测量、有序向量空间分离、定性概率与 Stieltjes 表示。本文与具体路径域相关的结果包括：精确路径纤维；由端点和总通行在一维中重构方向；在非加法封闭的历史域上进行可加定价；消失通行与保留目标达成之间的目标边界；以及即使只涉及六条历史，其系数重复次数也可能无界、但仍可由实际路径实现的消去证书。减烟是贯穿全文的方案环境；形式结果适用于具有相同通行语义的一维目标相对调整问题。

全文安排如下。第~\ref{sec:recovery}--\ref{sec:environment} 节定义经济对象与路径域。第~\ref{sec:universal} 节构造通用账户、精确纤维与一维压缩。第~\ref{sec:ordinal}--\ref{sec:stieltjes} 节建立可加支撑赋值以及正则测度化细化。第~\ref{sec:pha} 节给出端点--总通行--目标达成的解释和方案转移限制。第~\ref{sec:implications} 节讨论匹配设计与有限数据预测。附录集中给出纤维定理证明、分离论证、正则性构造及消去复杂度结果。

\section{目标相对的恢复历史}\label{sec:recovery}

\subsection{目标相对状态与四类有向运动}

令 $x_t$ 表示相对于一个预先给定目标测量的一维状态，并规定恰在目标处 $x_t=0$。在吸烟应用中，可以取 $x_t=A_t-A_t^*$，其中 $A_t$ 是成瘾存量、$A_t^*$ 是当前治疗目标；也可以采用任何保留目标位置和状态排序的单调有界归一化。本文把 $x_t$ 作为原始对象；具体应用另行给出其运动定律。

当 $x>0$ 时，向 $0$ 移动表示从目标上方进行减量，而远离 $0$ 表示复发加剧；当 $x<0$ 时，远离 $0$ 表示进一步低于目标或出现失稳，而向 $0$ 移动则表示回到计划治疗水平。由此定义四个模块 $T_R$、$A_R$、$A_L$ 和 $T_L$，每个模块都可以具有自身的价值测度。

\subsection{共同记账与个体特定赋值}

一份恢复日记首先生成一个有向通行账户。账户记录各模块跨越的目标相对状态水平及其重复次数。随后，原始条件弱序在把非通行后果视为固定背景变量的条件下，对各模块的通行剖面赋予主观价值。对于同一个个体，在整个比较域内使用同一套内部一致的赋值体系；对于另一个个体，相同账户可以得到不同的赋值。

这里的条件化明确了通行次序比较时保持不变的对象。在给定运动定律下，绕行与直接减量可能带来不同的消费、健康或调整后果，因此不必要求二者在这些后果字面完全相同时还能同时可行。条件次序所隔离的是对通行账户的一个部分评价。在人为设计或自然控制的菜单中，如果被比较的一对方案具有相同的非通行后果，则观察到的选择会直接揭示这一通行次序；若这些后果不同，则具体应用需要把通行次序与其他后果模型结合起来。与程序效用和序列效用类似 \citep{LoewensteinPrelec1993,FreyBenzStutzer2004,FreyStutzer2005}，这一条件对象可以具有经济含义，而无需在每个应用中都能把通行与其他后果分别操纵。

支撑通行赋值属于个体内部对象；若要进行个体间比较，还需要额外的归一化。

\begin{figure}[pos=htbp,align=\centering]
 \centering
 \includegraphics[width=0.98\textwidth]{Figure_1_recovery_accounting.png}
 \caption{\textbf{目标相对恢复历史的有向记账。} 图~(a) 展示一条依次接近、跨过、离开并再次到达目标的路径。图~(b) 把非静止的单调片段分配到 $T_L$、$T_R$、$A_L$ 和 $A_R$ 四个模块。图~(c) 给出目标界面的记账约定。一次连续跨越在物理上是一段运动，但在账户中分为两个有向通行片段：进入目标的趋向目标 片段包含目标并记录一次目标达成，而离开目标的远离目标 片段不包含目标处的质量。}
 \label{fig:accounting}
\end{figure}

\subsection{目标界面}

只要一个单调片段跨越目标，通行账户就在 $0$ 处将其拆分。进入目标的趋向目标 片段以目标为终点并包含目标点；离开目标的远离目标 片段从目标出发但排除目标点。因此，一次连续跨越恰好记录一次目标达成。尽管如此，这两个有向片段由于属于不同模块，其穿孔部分仍可具有不同价值。

例如，从 $x=2$ 连续移动到 $x=-2$ 在物理上是一次减量过程。在账户中，它由两部分构成：先是从右侧的 $2$ 向 $0$ 的趋向目标 通行，其中包含一次右侧目标达成；随后是从 $0$ 向 $-2$ 的左侧 远离目标 通行，且不计入目标处质量。这一记账方式把“到达目标”的价值与“越过目标继续前进”的价值分离开来。

\subsection{三种历史实验}

固定右侧距离 $0<a<b<c$。考虑三条具有共同初始距离 $c$ 和终止距离 $b$ 的历史：
\[
 \gamma^D:c\longrightarrow b,
 \qquad
 \gamma^C:c\longrightarrow a\longrightarrow b,
 \qquad
 \gamma^A:c\longrightarrow0\longrightarrow b.
\]
直接历史 $\gamma^D$ 记录单调减量；绕行历史 $\gamma^C$ 先进一步改善至 $a$，随后复发回到 $b$；触及目标的历史 $\gamma^A$ 先到达目标，再返回 $b$。三条历史具有相同的有序端点，因此它们之间的差异能够隔离内部通行价值。

对于任意归一化、正则且相容的支撑赋值 $\Lambda$，下文的表示推出
\[
\begin{aligned}
 U_\Lambda(\gamma^C)-U_\Lambda(\gamma^D)
  &=2\kappa_{R,\Lambda}((a,b]),\\
 U_\Lambda(\gamma^A)-U_\Lambda(\gamma^D)
  &=2\kappa_{R,\Lambda}((0,b])+\alpha_{R,\Lambda}.
\end{aligned}
\]
在一个固定支撑赋值下，第一个差分隔离了目标之外的“减量--复发”绕行价值；第二个差分则隔离了额外的近目标通行对比以及目标达成价值。在整个支撑族上，这些基数分量仅被集合识别。第~\ref{sec:pha} 节和第~\ref{sec:implications} 节把同一几何转化为序数转移限制和有限数据锐界。仅依赖端点的子类预测 $\gamma^D\sim\gamma^C\sim\gamma^A$；三者之间任何严格排名都会否定“目标相对终止状态充分性”。

\subsection{匹配持续时间与流量暴露}\label{sec:matched}

即使一些熟悉的结果汇总指标已经被匹配，方案的通行记录仍可以不同。令一个\emph{几何方案}指定有序的有限步状态路径；具体应用还可以为该路径提供一个可行的计时遍历。下面的结果保持每个计时遍历本身不变，只把共同端点处的停留用作匹配工具。

\begin{proposition}[固定遍历下的匹配方案设计]\label{prop:matched-design}
固定 $0\le b<c$，并给定有限多个包含在 $[0,c]$ 中的有限步几何方案 $\pi^1,\ldots,\pi^K$，每个方案都从 $c$ 出发并终止于 $b$。对每个 $k$，令 $z_k:[0,\tau_k]\to[0,c]$（$\tau_k>0$）是方案 $\pi^k$ 在应用中可行的一个计时遍历。假设可行时间安排类对在每个 $z_k$ 两端任意加入停留封闭：对任意 $\delta,\beta\ge0$，先在 $c$ 停留 $\delta$，随后不重参数化地沿 $z_k$ 运行，最后在 $b$ 停留 $\beta$，仍然可行。令连续函数 $g:[0,c]\to\R$ 满足 $g(c)\ne g(b)$。则存在有限的 $\overline T$，使得对每个 $T>\overline T$，都存在由 $z_k$ 仅通过加入端点停留而得到的可行时间安排 $\widetilde\pi^k:[0,T]\to[0,c]$，并且
\[
 \int_0^T g\bigl(\widetilde\pi^k(t)\bigr)\,\dd t
\]
对所有 $k$ 都相同。每个方案的非静止遍历保持不变，因此其通行账户也保持不变。
\end{proposition}

\begin{proof}
令
\[
 m_k=\int_0^{\tau_k} g(z_k(t))\,\dd t,
 \qquad
 \Delta=g(c)-g(b),
 \qquad
 a_k=\frac{\tau_k g(b)-m_k}{\Delta}.
\]
选择
\[
 \overline T
 =2\max_{1\le k\le K}\{0,-a_k,\tau_k+a_k\}.
\]
对于任意 $T>\overline T$，定义初始与终止停留时间
\[
 \delta_k=\frac{T}{2}+a_k,
 \qquad
 \beta_k=\frac{T}{2}-\tau_k-a_k.
\]
根据 $\overline T$ 的定义，二者均严格为正，且 $\delta_k+\tau_k+\beta_k=T$。第 $k$ 个时间安排先在 $c$ 停留 $\delta_k$，然后不作重参数化地沿原计时遍历 $z_k$ 运行，最后在 $b$ 停留 $\beta_k$。其累计流量为
\[
 \delta_k g(c)+m_k+\beta_k g(b)
 =\frac{T}{2}\bigl(g(c)+g(b)\bigr),
\]
其中最后一个等式使用了 $a_k\Delta=\tau_k g(b)-m_k$。因此，所有时间安排具有相同的持续时间和累计流量。我们只增加了端点停留，所以根据假设可行性得到保留，同时通行账户不变。
\end{proof}

取 $g(x)=x$ 即可匹配累计目标相对暴露。特别地，上一小节中直接历史、目标外绕行历史以及触及目标历史的任意应用可行计时版本，都可以嵌入一个足够长的时间安排中，在不改变遍历速度的情况下实现共同端点、共同持续时间和共同累计暴露。条件 $g(c)\ne g(b)$ 是非退化条件：在两个端点之间重新分配停留时间会改变累计流量，却不改变持续时间和通行记录。作为特例，如果遍历时间本身可以自由缩放，则时间跨度限制消失：对任意给定的 $T>0$，可以先给有限个遍历安排足够短的共同遍历时间，然后再应用端点停留构造。该设计可以精确匹配一个选定的、能够区分端点的连续流量统计量；若还要匹配其他后果，则需要更多设计变量或一个明确的应用模型。

\subsection{比较域与可行选择}

表示域包含用来比较通行记录并识别其代数限制的有限步目标相对历史。每一个治疗问题都有自己的可行菜单，而它可能只是这一比较域中的一个很小子集。因此，循环是否可重复、是否存在反向机会以及是否允许触及目标，都是应用特定的可行性特征。

受控比较使通行次序具有经验解释。命题~\ref{prop:matched-design} 表明：一旦给定应用上可行的遍历，只要共同时间跨度足够长，就能仅通过加入端点停留，精确匹配持续时间和一个能够区分端点的累计流量后果，而无需改变遍历速度。其他后果可以通过更丰富的设计进行匹配、单独测量，或由提供可行菜单的具体应用模型处理。

\section{与文献的关系}\label{sec:literature}

\paragraph{戒烟、理性成瘾与目标相对恢复。}
戒烟构成全文贯穿的方案环境。\citet{BeckerMurphy1988} 用成瘾存量和相邻互补性刻画稳定的前瞻偏好，\citet{BeckerGrossmanMurphy1994} 则把这一结构用于香烟需求。后续研究讨论变化的短视程度 \citep{OrphanidesZervos1998}、现在偏误与自我控制需求 \citep{GruberKoszegi2001}、戒断与戒烟决策 \citep{SuranovicGoldfarbLeonard1999}、线索触发选择 \citep{BernheimRangel2004}，以及强迫性消费中的反复 \citep{Winston1980}。本文的通行次序隔离了一个互补对象：在给定上述模型所表示的其他后果条件下，个体如何对有向调整记录本身进行排序。

预先给定但不具强制约束力的目标，可以理解为治疗参照点或分阶段减量平台。目标可以支持自我调节 \citep{Hsiaw2013,KochNafziger2011}，承诺契约也已经直接用于戒烟研究 \citep{GineKarlanZinman2010}。个体可以未达到目标、越过目标或远离目标。内生目标形成、愿望水平修订和承诺机制设计属于应用层面的进一步扩展。

\paragraph{参照依赖与目标型状态效用。}
参照依赖模型相对于一个参照点评价结果，并允许参照点两侧具有不对称的敏感度 \citep{TverskyKahneman1991,KoszegiRabin2006}。本文表示中最接近这类模型的是仅依赖端点的子类。在任意固定的正则相容支撑赋值下，当 $\kappa_L=\kappa_R=0$ 且 $\alpha_L=\alpha_R=0$ 时，通行价值退化为
\[
 U_\Lambda(\gamma)=P_\Lambda(\gamma(1))-P_\Lambda(\gamma(0)).
\]
在初始状态固定时，每个支撑赋值于是只依赖以目标为中心的终止状态。一般表示在支撑赋值层面额外加入两种路径价值来源：由有向跨越次数产生的总通行调整，以及由实际目标达成产生的边界贡献。本文维持目标预先给定的环境，把基于预期形成参照点的问题留给具体应用。

\paragraph{序列、经历与程序价值。}
关于序列和延展经历的偏好可见于 \citet{LoewensteinPrelec1993} 与 \citet{VareyKahneman1992}。程序效用允许结果的生成方式本身进入福利 \citep{FreyBenzStutzer2004,FreyStutzer2005}；\citet{MinardiSavochkin2024} 则公理化了过去消费的可记忆效应。本文的通行分量保留一侧、方向、状态水平与重复次数，同时忽略日历时间、持续时间、新近性、适应以及彼此独立的闭合经历之间的顺序。

形如 $g(c,A)=\int_c^A q(s)\,\dd s$ 的参数积分为一次单调调整所跨越的状态水平赋值。有向通行测度把这一思路以序数、非参数方式扩展到重复的非单调历史、目标两侧、两个方向以及目标达成原子。

\paragraph{可加测量与受限定义域。}
从有限消去条件到可加表示的路线是经典的；参见 \citet{Scott1964}、\citet{KrantzEtAl1971} 和 \citet{Wakker1989}。\citet{KraftPrattSeidenberg1959} 在有限定性概率中发现的障碍，与本文使用平衡的有限证书而不只依赖成对一致性的做法非常接近。\citet{Fishburn2001} 研究有限乘积域上所需的不同比较对数量，同时允许证书中出现整数重复次数。本文的算术构造固定六个账本列，并让一个路径可实现族中的原始系数重复次数发散。

部分定义的运算是广延测量中的核心问题。\citet{Luce1967} 在不交并仅部分定义时得到有限可加概率表示，\citet{Mackenzie2021} 则从原子聚集条件推导相关假设。本文通过一个通用通行群对可行路径拼接作代数完备化，但只对实际比较所生成的揭示锥进行排序。受限域上的可加表示见 \citet{ChateauneufWakker1993}、\citet{Kopylov2007} 和 \citet{QinRommeswinkel2022}。允许形式差分的次序保持偏序，使实数定价步骤更接近多效用表示，而不是全局 H\"older 标量化；参见 \citet{Ok2002}、\citet{DubraMaccheroniOk2004} 和 \citet{EvrenOk2011}。

\paragraph{连续性、原子、路径与原始平衡。}
在定性概率中，单调连续性是从有限可加性走向可数可加性的经典路线 \citep{Villegas1964,ChateauneufJaffray1984,Chateauneuf1985}；相关连续性限制也被用于正则化先验族 \citep{ChateauneufMaccheroniMarinacciTallon2005}。本文的连续性条件由路径几何生成：只有几何交集为空的非目标穿孔区间被视为可忽略，而到达目标的序列可以保留一个单点残差。

忽略顺序而保留重复次数的汇总与 Parikh 向量有关 \citep{Parikh1966}，而路径签名则保留非交换的顺序信息 \citep{Chen1954,HamblyLyons2010}。纤维定理把 Euler 迹变换 \citep{Abrham1980,FleischnerEtAl1992} 专门化到带标签的一维路径，并固定端点不平衡与目标达成标签。总水平跨越计数有经典分析传统 \citep{Lochowski2017}；在本文中，它们被用于偏好表示以及端点--总通行分解。

回路、共形不可分解性和 Graver 基都是标准概念 \citep{Graver1975}。相关复杂度结果包括 \citet{BersteinOnn2009} 与 \citet{TatakisThoma2015}。本文的原始元素本身就是一个回路，对每个参数值都只有六列，其增长来自可达路径账户中的算术系数。总通行调整是一个静态、对反转保持偶对称的分量。进入--退出不可逆性 \citep{Dixit1989}、动态迟滞 \citep{MielkeRoubicek2015}，以及博弈和排序中的势--循环分解 \citep{CandoganEtAl2011,JiangEtAl2011}，提供了相邻但不同的结构。

\section{恢复路径与有向账本}\label{sec:environment}

尽管形式结果适用于任何一维目标相对调整，全文的贯穿解释仍是相对于分阶段治疗目标的成瘾存量。形式路径是一份恢复日记；可达性问题询问：给定的阈值跨越与端点是否能够共同属于同一条连续历史。

\subsection{路径与模块}

令
\[
 X=[\underline x,\overline x]\subset\R,
 \qquad \underline x<0<\overline x,
\]
并固定目标为 $0$。记 $R_L=-\underline x$、$R_R=\overline x$。所谓\emph{有限步路径}，是连续映射 $\gamma:[0,1]\to X$，使得存在 $[0,1]$ 的某个有限分割，并且 $\gamma$ 在每个分割单元上单调。允许存在常值片段。

四个通行模块为
\[
 \M=\{T_L,T_R,A_L,A_R\}.
\]
符号 $T$ 表示朝向目标的运动，$A$ 表示远离目标的运动；下标表示所在一侧。对 $m\in\M$，令 $R_m$ 表示相应一侧上的最大目标距离。

对 $0\le a<b\le R_m$，令 $p_m(a,b)$ 表示跨越距离区间 $[a,b]$ 的典范单调通行。趋向目标 通行从距离 $b$ 移向 $a$；远离目标 通行从 $a$ 移向 $b$。端点约定是不对称的：
\begin{itemize}
 \item 趋向目标 通行 $p_{T_s}(0,b)$ 记录目标达成，因此覆盖 $[0,b]$；
 \item 远离目标 通行 $p_{A_s}(0,b)$ 记录经 $(0,b]$ 离开目标的过程，并且在 $0$ 处不另计单点通行。
\end{itemize}

\subsection{通行剖面}

每个单调片段在其所属模块中贡献一个阶梯函数。对所有片段求和得到有向通行剖面
\[
 h^m_\gamma:[0,R_m]\to\N_0,
 \qquad m\in\M.
\]
对于穿孔区间 $(a,b]$，$h^m_\gamma(r)$ 计数模块 $m$ 中有多少片段跨越状态水平 $r$。在目标处，$h^{T_s}_\gamma(0)$ 是从侧 $s$ 到达目标的次数，而 $h^{A_s}_\gamma(0)=0$。具体账本为
\[
 \ell(\gamma)
 =\bigl(h^{T_L}_\gamma,h^{T_R}_\gamma,h^{A_L}_\gamma,h^{A_R}_\gamma\bigr)
 \in\Led.
\]
形式上，$\Led$ 是四个整数阶梯函数群的直积。趋向目标 一侧的群由 $\1_{[0,b]}$ 与满足 $0<a<b$ 的 $\1_{(a,b]}$ 生成，而 远离目标 一侧的群由满足 $0\le a<b$ 的 $\1_{(a,b]}$ 生成。实际账本构成该群中非负且可达的子集。
账户不记录各片段的具体先后顺序，也不记录遍历这些片段所花费的时间，但保留重复次数。

一个有限账本在其断点所出现的状态上诱导出一个有限的带标签有向多重图。每一次对相邻状态区间的单位跨越都对应一条按所在一侧和方向标记的弧。

\begin{proposition}[记录通行账户的可行性]\label{prop:reachability}
令 $L$ 为有限的四模块阶梯账本，并令 $s,t\in X$ 为给定的有序端点。
\begin{enumerate}
 \item 若 $L\ne0$，则 $L$ 能由一条从 $s$ 到 $t$ 的有限步路径生成，当且仅当其非孤立的带标签多重图连通，并且
 \[
  \deg^+(v)-\deg^-(v)
  =\1_{\{v=s\}}-\1_{\{v=t\}}
 \]
 对每个顶点 $v$ 都成立；此外，当 $s=t$ 时，共同端点 $s$ 必须属于非孤立支撑。
 \item 若 $L=0$，则存在一条从 $s$ 到 $t$ 的路径生成 $L$，当且仅当 $s=t$；此时实现路径在弱时间变换意义下都是常值路径。
\end{enumerate}
对于满足第~1 项条件的非零账本，带标签多重图的一条 Euler 迹给出一条实现路径。
\end{proposition}

\begin{proof}[证明思路]
一条实现路径必须按照记录的重复次数遍历每条带标签弧，因此强制要求非孤立支撑连通并满足所显示的端点不平衡条件。反过来，把重复次数替换成相应数量的平行弧以后，这些条件恰好就是有向 Euler 迹存在的判据。沿 Euler 迹读取顶点并在相邻状态之间插值得到一条实现用的有限步路径；零账本只能由常值路径实现。完整证明见附录~\ref{app:universal}。
\end{proof}

命题~\ref{prop:reachability} 是基本的定义域限制。账本的形式和未必满足其连通性或端点平衡条件，因此实际路径域并不对加法封闭。

\subsection{偏好}

原始偏好 $\succeq$ 是所有有限步路径集合 $\Pth$ 上的条件弱序，其中非通行后果保持固定。状态 $x$ 处的常值路径记为 $c_x$。除非某个结果明确把注意力限制在有限菜单上，否则 $\Pth$ 表示 $X$ 上完整的有限步路径域。特别地，它包含所有常值路径、反向路径、端点相容的拼接、典范通行、典范目标外往返、触及目标的循环，以及把任意典范区间细分成任意给定有限个非空典范区间所得的路径。各项假设承担四种不同角色：弱序和 BPCC 约束实际比较的可加一致性；PVBC 和 PROC 分别约束实数完备化与穿孔连续性；典范区间丰富性是定义域条件；方向单调性则给目标相对模块提供经济学方向。

\begin{assumption}[面向目标的方向单调性]\label{ass:directional}
对任意非空典范区间 $I$，每一个趋向目标 通行 $p_{T_s}(I)$ 都严格优于目标处的常值路径 $c_0$，而 $c_0$ 又严格优于每一个远离目标 通行 $p_{A_s}(I)$。在 BPCC 下，所有常值路径具有相同的零通行账户，因此彼此无差异。
\end{assumption}

方向单调性赋予预先给定的目标一个行为方向：在维持的背景后果条件下，局部朝向目标的运动优于停留，而局部远离目标的运动更差。在正则表示中，它推出非负的目标达成残差；但达成残差是否严格为正仍是偏好的额外性质。BPCC 则从相同账本推出纤维内无差异。

\section{通用通行账户与通行等价日记}\label{sec:universal}

第一项任务是把通行的共同描述与个体赋值分离开来。即使已经指定运动发生在哪里、位于目标哪一侧、方向如何以及重复多少次，路径的效用仍未确定。通用定理识别出所有允许的赋值都必须经过的可加信息空间；纤维定理随后精确刻画：何时两份有序日记会向这一赋值类中的每一个赋值提供完全相同的输入。

\subsection{通用可加载体}

对每个 $m\in\M$ 以及 $0\le a<b\le R_m$，引入一个形式生成元 $e_m(a,b)$。令 $\mathbb Z^{(\mathcal E)}$ 为这些生成元上的自由阿贝尔群，并对下列细分关系取商：
\[
 e_m(a,c)=e_m(a,b)+e_m(b,c),
 \qquad 0\le a<b<c\le R_m.
\]
所得群记为 $\mathfrak L$。在路径的转折点和目标跨越点处切分，可定义
\[
 q:\Pth\to\mathfrak L
\]
即把相应生成元的等价类相加。细分关系保证 $q$ 与所选切分方式无关。

弱时间变换是与一个从 $[0,1]$ 到自身的连续、非减、满射映射作复合。它可以改变速度并插入停留，但不能颠倒访问状态的先后顺序。

对于阿贝尔群 $G$，令 $\operatorname{Val}_G(\Pth)$ 表示所有映射 $\Phi:\Pth\to G$ 的集合，这些映射对弱时间变换不变，并且对每一次端点相容的拼接可加。

\begin{theorem}[可加赋值的通用通行账户]\label{thm:universal}
令 $G$ 为任意阿贝尔群。映射 $\Phi:\Pth\to G$ 对弱时间变换不变且对端点相容拼接可加，当且仅当存在唯一的同态 $\bar\Phi:\mathfrak L\to G$ 使得
\[
 \Phi=\bar\Phi\circ q.
\]
等价地，
\[
 \operatorname{Hom}_{\mathbf{Ab}}(\mathfrak L,G)
 \cong
 \operatorname{Val}_G(\Pth).
\]
此外，$\mathfrak L$ 与四模块整数阶梯账本群典范同构，并且在该同构下 $q$ 恰好变为 $\ell$。
\end{theorem}

由于该定理本身不含具体赋值，它同时允许个体特定和模块特定的权重。对于个体 $i$，任何允许的通行赋值都具有形式
\[
 \Phi_i(\gamma)=\bar\Phi_i\bigl(q(\gamma)\bigr),
\]
其中 $\bar\Phi_i$ 是个体特定同态。共同的是账户 $q(\gamma)$；不同的是对该账户的赋值。

常值路径的零通行价值由 $c_x*c_x=c_x$ 与可加性推出。状态水平、持续时间、健康、消费等非通行后果在条件通行排序中保持固定。

定理~\ref{thm:universal} 使用完整定义域 $\Pth$，其中包含每个典范生成元以及细分所需的所有拼接。受限菜单则需要各自的丰富性条件。

\subsection{通行账户的精确纤维}

令 $\approx_p$ 为下列操作生成的等价关系：
\begin{enumerate}
 \item 弱时间变换，包括插入或删除停留；
 \item 把所有常值路径认作零通行类；
 \item 交换以同一状态为基点的两个闭合游程；
 \item 对闭合路径作循环旋转，即对同一闭合日记选择不同的记录起点。
\end{enumerate}
每一种操作都保持有向账户不变。逆命题更微妙：需要证明账本系数相同的两条路径可以通过一系列操作连接，并且每个中间对象仍是可行的单一路径，同时目标达成标签始终保持不变。

\begin{theorem}[通行等价的恢复日记]\label{thm:fiber}
对于固定目标区间上的有限步路径，
\[
 q(\gamma)=q(\eta)
 \quad\Longleftrightarrow\quad
 \gamma\approx_p\eta.
\]
\end{theorem}

证明分三步使用几何结构。首先，把有限多条连续路径约化为同一个带标签路径图上的顶点词；其次，利用流量不平衡恢复每个开放实现的有序端点；最后，通过带根的叶节点删除，把每个实现写成一条约化核心路径加上一组闭合叶游程。闭合游程交换和插入点不变性使这些多重集能够对齐。只有对非零闭合账本选择共同基点时才需要循环旋转。若删除的叶节点就是目标，则恢复的游程恰好包含一个以目标为终点的趋向目标 单元和一个从目标出发的远离目标 单元，因此目标达成约定得到保留。完整论证见附录~\ref{app:fiber}。图~\ref{fig:fiber} 展示同基点闭合游程交换这一生成操作。

\begin{figure}[pos=htbp,align=\centering]
 \centering
 \includegraphics[width=0.98\textwidth]{Figure_2_path_fiber.pdf}
 \caption{\textbf{路径纤维内部的闭合游程交换。} 两份有序日记共享同一开放骨架，并包含以目标为共同基点的两个相同闭合游程，但先后顺序相反。交换这些同基点游程会保持每一个穿孔有向剖面以及各侧的目标达成次数。因此 $q(\gamma)=q(\eta)$ 且 $\gamma\approx_p\eta$。该图展示了路径纤维等价关系的一个生成操作。}
 \label{fig:fiber}
\end{figure}

每一个可加通行赋值都会给账本等价的历史赋予相同通行价值。上述操作穷尽了有向通行账户所丢弃的信息；在更大的应用模型中，非通行后果仍可能区分这些历史。

该结果利用有限的一维路径几何。经典的片段反转与分离--吸收结果允许更广泛的变换 \citep{Abrham1980,FleischnerEtAl1992}；本文定理则只使用始终停留在实际路径域中的账户保持操作，给出精确的纤维描述。

\subsection{由总通行与端点得到的一维压缩}\label{sec:compression}

对于侧 $s\in\{L,R\}$ 和穿孔距离 $r>0$，令 $N_s^T(r;\gamma)$ 与 $N_s^A(r;\gamma)$ 分别表示 趋向目标 与 远离目标 的跨越次数，并定义
\[
 D_s^\gamma=N_s^T-N_s^A,
 \qquad
 V_s^\gamma=N_s^T+N_s^A.
\]
当 $x$ 位于侧 $s$ 且距目标至少为 $r$ 时，令 $\chi_s(x,r)=1$；否则令其为零。

\begin{lemma}[净通行守恒]\label{lem:net}
对任意有限步路径和任意穿孔距离水平，
\[
 D_s^\gamma(r)
 =\chi_s(\gamma(0),r)-\chi_s(\gamma(1),r).
\]
因此，净有向通行完全由有序端点决定。
\end{lemma}

\begin{proposition}[由总通行与有序端点无损重构]\label{prop:gross-endpoints}
在有限步一维路径域上，记录
\[
 \bigl(V_L^\gamma,V_R^\gamma,\gamma(0),\gamma(1)\bigr)
\]
即可重构完整的有向通行账户 $q(\gamma)$。因此，在 BPCC 下，该记录相同意味着原始偏好上的无差异。
\end{proposition}

\begin{proof}
引理~\ref{lem:net} 由有序端点恢复 $D_s$，因此
\[
 N_s^T=(V_s+D_s)/2,
 \qquad
 N_s^A=(V_s-D_s)/2.
\]
对于有限步路径，$N_s^T$ 在零点附近足够小的穿孔邻域上为常数。其稳定的近目标值恰好等于实际从侧 $s$ 到达目标的趋向目标 通行次数。因此，穿孔有向剖面与目标轨迹都得到恢复，这正是完整账本。
\end{proof}

在严格方向单调性下，仅有总通行并不充分：同一穿孔区间上的典范 趋向目标 通行与其反向 远离目标 通行具有相同的总剖面，却分别位于停留的不同偏好方向。有序端点提供了缺失的方向不平衡。这一重构结果依赖于一维有限步几何。

\section{全局一致性与支撑通行赋值}\label{sec:ordinal}

通用账户只描述通行而不对其赋值。下一步的问题是：一个个体对实际历史的比较，能否在账户空间上支撑一致的可加赋值。BPCC 使揭示出来的可加次序在实际历史上保持精确，同时保留未揭示形式差分上的偏序性。PVBC 则刻画何时阿基米德完备化能够保持这些实际历史比较。由此得到的表示通常是一族归一化的实数支撑赋值。

\subsection{典范有序包络}

令 $\succeq$ 为 $\Pth$ 上的一个弱序。定义比较锥
\[
 \mathcal C_{\succeq}
 =\left\{
  \sum_{i=1}^n\bigl(q(\gamma_i)-q(\eta_i)\bigr):
  \gamma_i\succeq\eta_i
 \right\}\subseteq\mathfrak L,
\]
对称子群
\[
 \mathcal N_{\succeq}
 =\mathcal C_{\succeq}\cap(-\mathcal C_{\succeq}),
\]
以及商群
\[
 G_{\succeq}=\mathfrak L/\mathcal N_{\succeq}.
\]
比较锥的像在 $G_{\succeq}$ 上定义一个偏序。记 $J_{\succeq}=\pi_{\succeq}\circ q$。

$\mathcal C_{\succeq}$ 中的有限和就是一致性证书：它们把实际历史比较所揭示的账户差分聚合起来，即使最终得到的形式和本身位于可行历史域之外。

\begin{definition}[平衡通行循环一致性（BPCC）]
若不存在有限列表 $\gamma_i\succeq\eta_i$，其中至少有一个严格比较，并满足
\[
 \sum_i q(\gamma_i)=\sum_i q(\eta_i),
\]
则称 BPCC 成立。
\end{definition}

\begin{proposition}[典范有序包络与 BPCC 的精确性]\label{prop:envelope}
对任意弱序，$G_{\succeq}$ 都是一个偏序阿贝尔群，且 $J_{\succeq}$ 具有单调性、弱时间变换不变性并对拼接可加。BPCC 成立，当且仅当
\[
 \gamma\succeq\eta
 \quad\Longleftrightarrow\quad
 J_{\succeq}(\gamma)\ge J_{\succeq}(\eta).
\]
任意取值于有序阿贝尔群、且单调可行的赋值，都通过一个正同态唯一地因子分解于 $J_{\succeq}$。
\end{proposition}

该命题赋予 BPCC 三个精确作用：它排除平衡的严格改进循环；使典范映射在实际历史上精确；并使典范商在所有单调可加赋值中具有通用性。特别地，只有一次比较时得到
\[
 q(\gamma)=q(\eta)
 \quad\Longrightarrow\quad
 \gamma\sim\eta.
\]
因此，通行纤维上的偏好无差异是由 BPCC 推出的，而不是作为原始公理直接加入。

\subsection{实际历史次序上的成对实数分离}

在实数分离步骤中，记
\[
 d(\gamma,\eta)=q(\gamma)-q(\eta),
 \qquad
 E=\mathfrak L\otimes_{\Z}\R,
\]
并定义实数揭示锥与无差异空间
\[
 C=\cone\{d(\gamma,\eta):\gamma\succeq\eta\},
 \qquad
 I=\spanr\{d(\gamma,\eta):\gamma\sim\eta\}.
\]
任意有限个通行向量都位于某个共同有限网格上，因此实锥证书可以有理化并通过清除分母转化为整数证书。

\begin{proposition}[精确实数化揭示锥]\label{prop:realcone}
在 BPCC 下，
\[
 C\cap(-C)=I,
\]
并且对于实际历史，
\[
 d(\gamma,\eta)\in C
 \quad\Longleftrightarrow\quad
 \gamma\succeq\eta.
\]
\end{proposition}

转到 $V=E/I$ 与 $K=C/I$。对于一个非空典范区间 $J$，定义其有利账户差分为
\[
 x_m(J)=
 \begin{cases}
  d(p_m(J),c_0),&m\in\{T_L,T_R\},\\
  d(c_0,p_m(J)),&m\in\{A_L,A_R\}.
 \end{cases}
\]
在假设~\ref{ass:directional} 下，在每个模块中选取一个全范围区间 $J_m$，并令
\[
 u=\sum_{m\in\M}x_m(J_m),
 \qquad e=u+I\in V.
\]
细分关系使 $e$ 成为序单位：对每个 $v\in V$，都存在有限的 $r>0$ 使其位于 $-re$ 与 $re$ 之间。任意两个全支撑基准都会生成同一个完备化，因为它们彼此都能被有限正倍数支配。

定义阿基米德闭包
\[
 K^{\mathrm a}
 =\left\{v\in V:\ v+n^{-1}e\in K\text{ 对每个 }n\in\N\right\},
\]
其线性部分 $N=K^{\mathrm a}\cap(-K^{\mathrm a})$，以及阿基米德化的序单位空间
\[
 \widehat V=V/N,
 \qquad
 \widehat K=K^{\mathrm a}/N,
 \qquad
 \widehat e=e+N.
\]
令 $\rho:V\to\widehat V$ 为商映射，并记
\[
 \widehat d_{\gamma\eta}=\rho(d(\gamma,\eta)+I).
\]

\begin{definition}[成对消失基准一致性（PVBC）]\label{def:pvbc}
若对每一对实际历史 $\gamma,\eta$，都有
\[
 d(\gamma,\eta)+I\in K^{\mathrm a}
 \quad\Longrightarrow\quad
 \gamma\succeq\eta.
\]
等价地，若 $\eta\succ\gamma$，则存在有限重复次数 $n=n(\eta,\gamma)$，使得一个固定全支撑基准不再支持“用 $n$ 份 $\gamma$ 补偿 $n$ 份 $\eta$”这一派生比较。不同严格比较之间不要求存在共同的重复次数上界。
\end{definition}

对小明而言，PVBC 的含义是：如果两条戒烟历史之间存在真实劣势，那么当同一劣势重复足够多次时，它最终会超过一个固定恢复基准。所需重复次数依赖于具体比较对，并且可以任意大。该条件排除了实际历史上的无穷小严格排序，同时保留未揭示形式组合的偏序性，也不限制闭合游程的取值。

归一化支撑赋值是一个线性泛函 $\varphi:\widehat V\to\R$，满足
\[
 \varphi(\widehat K)\subseteq\R_+,
 \qquad
 \varphi(\widehat e)=1.
\]
令 $\mathcal S$ 表示所有此类赋值的集合，并定义
\[
 U_\varphi(\gamma)=\varphi\!\left(\rho(q(\gamma)+I)\right).
\]

\begin{theorem}[实际历史上的成对实数分离]\label{thm:A}
假设 BPCC 与严格方向单调性成立。则 PVBC 成立，当且仅当归一化支撑族 $\mathcal S$ 精确表示实际路径次序：
\[
 \gamma\succeq\eta
 \quad\Longleftrightarrow\quad
 U_\varphi(\gamma)\ge U_\varphi(\eta)
 \quad\text{对每个 }\varphi\in\mathcal S.
\]
族 $\mathcal S$ 非空、凸，并且在 $\widehat V$ 上逐点收敛拓扑下为紧集。若 $\eta\succ\gamma$，则每个归一化支撑赋值都弱排序 $\eta$ 高于 $\gamma$，且至少有一个支撑赋值给出严格排序。
\end{theorem}

\begin{proof}[证明思路]
锥 $\widehat K$ 是闭的、尖的，并具有序单位 $\widehat e$。因此正分离给出
\[
 \widehat v\in\widehat K
 \quad\Longleftrightarrow\quad
 \varphi(\widehat v)\ge0\quad\forall\varphi\in\mathcal S.
\]
实际弱比较位于 $K\subseteq K^{\mathrm a}$。反过来，一致同意把一个实际差分放入 $K^{\mathrm a}$，PVBC 再把它还原成原始弱比较。逆向论证给出必要性。关于实锥、阿基米德化、分离与紧性的细节见附录~\ref{app:ordered}。
\end{proof}

不同的归一化支撑赋值可以在那些未被实际历史比较固定的形式方向上意见不一，但它们的一致同意仍恢复实际历史上的原始次序。因此，单个 $U_\varphi$ 可能把一个原始严格比较判为无差异；严格性只保证在支撑族中的至少某个成员上出现。下面的单点情形则是一个特殊情况：一个归一化赋值单独就忠实表示次序。

\begin{remark}[有限网格并不会自动保证 PVBC]\label{rem:finitegrid-pvbc}
有限物理网格使账户空间有限维，却不限制路径重复次数，因此它本身并不能排除无穷小严格比较。例如，在一个网格上只有一个有利的趋向目标 单元 $T$ 及其不利的反向 远离目标 单元 $A$，在字典序群 $\Z^2$ 中赋值
\[
 \Phi(T)=(1,0),
 \qquad
 \Phi(A)=(-1,1).
\]
按照通行账户的字典序值对实际历史排序，并给其他有利单元赋予第一坐标为正的向量。于是 $T\succ0\succ A$，且由于偏好由一个可加线性有序阿贝尔群诱导，BPCC 成立。闭合循环 $T+A$ 的值为 $(0,1)>_{\mathrm{lex}}0$，然而无论把该循环重复有限多少次，都无法超过任何第一坐标为正的全支撑基准。因此，在有限网格上 BPCC 与方向单调性可以同时成立而 PVBC 失败。有限经验设计则不同：尽管无限制的有限网格路径次序未必是阿基米德的，但有限多个已观察不等式会生成一个多面体补全问题。
\end{remark}

BPCC 是通行可加性背后的有限列表一致性条件。附录~\ref{app:primitive} 通过构造只含六个备选项但原始平衡证书所需重复次数无界的族，衡量这一条件的全局深度。这种有限证书逻辑与定性概率中的消去条件接近 \citep{KraftPrattSeidenberg1959,ChateauneufJaffray1984}，而支撑族分离表示则与多效用理论相联系 \citep{Ok2002,DubraMaccheroniOk2004,EvrenOk2011}。

\section{穿孔揭示次序连续性与目标达成}\label{sec:stieltjes}

实数支撑族在目标边界处仍留下一个连续性问题。穿过穿孔区间 $(a,b_n]$ 的典范通行在 $b_n\downarrow a$ 时从通行记录中消失；而穿过 $[0,b_n]$ 的到达目标通行始终包含目标达成事件。PROC 要求揭示比较对于第一类消失通行保持闭合。附录~\ref{app:regular} 构造相应的局部凸拓扑，并识别其连续正泛函。

\begin{assumption}[典范区间丰富性]\label{ass:richness}
对每个模块 $m$、每个非空典范区间 $I$ 以及每个整数 $q\ge1$，区间 $I$ 都能被划分为 $q$ 个非空典范区间，并且相应通行均属于 $\Pth$。各侧范围 $[0,R_m]$ 有限。
\end{assumption}

对每个模块 $m$，令 $\widehat x_m(J)=\rho(x_m(J)+I)\in\widehat K$ 表示穿过穿孔区间 $J$ 的有利典范通行。当 $I_n=(a,b_n]\downarrow\varnothing$ 且 $b_n\downarrow a$ 时，称
\[
 \sigma=(m,(I_n)_{n\ge1})
\]
为一个可容许穿孔收缩序列；其中 远离目标 模块允许 $a\ge0$，趋向目标 模块要求 $a>0$。到达目标的序列 $[0,b_n]\downarrow\{0\}$ 则构成一个单独的边界类别。

附录~\ref{app:regular} 把这种几何意义上的“消失”转化为一个局部凸拓扑 $\tau^\circ$。上述非目标穿孔区间收敛到零，而到达目标的序列留在边界类中。若归一化支撑赋值 $\varphi\in\mathcal S$ 对每个可容许穿孔收缩序列都满足
\[
 \varphi(\widehat x_m(I_n))\longrightarrow0,
\]
则称其为\emph{穿孔正则}。令 $\mathcal S_{\mathrm{reg}}$ 表示正则支撑赋值集合。附录~\ref{app:regular} 证明：它们恰好是拓扑 $\tau^\circ$ 下连续的归一化正泛函。

\begin{assumption}[穿孔揭示次序闭性（PROC）]\label{def:proc}
对任意一对实际历史 $\gamma,\eta$，若
\[
 \widehat d_{\gamma\eta}\in\overline{\widehat K}^{\,\tau^\circ},
\]
则
\[
 \gamma\succeq\eta.
\]
\end{assumption}

PROC 有一个成对的行为解释：若一系列揭示弱比较通过那些最终从记录中消失的非目标通行碎片趋近某一对实际历史，则极限处的实际历史对仍保持弱排序。对小明而言，像 $(10,10+\varepsilon_n]$、$\varepsilon_n\downarrow0$ 这样的区间会作为调整事件消失，但“到达治疗目标”仍作为边界事件保留下来。

\begin{theorem}[带目标达成的正则有向测度表示]\label{thm:B}
假设 BPCC、严格方向单调性、PVBC 以及假设~\ref{ass:richness} 成立。则 PROC 成立，当且仅当正则支撑族精确表示实际路径次序：
\[
 \gamma\succeq\eta
 \quad\Longleftrightarrow\quad
 U_\varphi(\gamma)\ge U_\varphi(\eta)
 \quad\text{对每个 }\varphi\in\mathcal S_{\mathrm{reg}}.
\]
在这些条件下，每个 $\varphi\in\mathcal S_{\mathrm{reg}}$ 都诱导出唯一的归一化四测度系统
\[
 \Lambda_\varphi=(\lambda_{T_L}^\varphi,\lambda_{T_R}^\varphi,
                   \lambda_{A_L}^\varphi,\lambda_{A_R}^\varphi),
\]
其中各分量都是相应侧范围上的有限非负 Borel 测度，并满足
\[
\begin{aligned}
 U_\varphi(\gamma)
 ={}&\int h^{T_L}_\gamma\dd\lambda_{T_L}^\varphi
   +\int h^{T_R}_\gamma\dd\lambda_{T_R}^\varphi\\
  &-\int h^{A_L}_\gamma\dd\lambda_{A_L}^\varphi
   -\int h^{A_R}_\gamma\dd\lambda_{A_R}^\varphi,
\end{aligned}
\]
且
\[
 \lambda_{A_s}^\varphi(\{0\})=0,
 \qquad
 \lambda_{T_s}^\varphi=\alpha_s^\varphi\delta_0+\lambda_{T_s}^{\varphi,\circ},
 \qquad
 \alpha_s^\varphi\ge0.
\]
令 $\mathfrak M_{\mathrm{reg}}(\succeq)$ 表示这些归一化测度系统的集合，则
\[
 \gamma\succeq\eta
 \quad\Longleftrightarrow\quad
 U_\Lambda(\gamma)\ge U_\Lambda(\eta)
 \quad\text{对每个 }\Lambda\in\mathfrak M_{\mathrm{reg}}(\succeq).
\]
该族可以包含多个归一化支撑赋值，并且不一定紧。严格方向单调性使每个典范有利通行在整个正则族上都取非负值，并且至少在一个成员上严格为正。
\end{theorem}

\begin{proof}[证明思路]
令 $K^\circ=\overline{\widehat K}^{\,\tau^\circ}$。若一个实际差分位于 $K^\circ$ 之外，则局部凸分离给出一个连续正泛函，并且该泛函在此差分上严格为负。用序单位归一化后得到正则状态，因为 $\tau^\circ$ 的连续对偶恰好就是穿孔正则对偶。因此，PROC 使正则状态足以进行精确分离；反方向由连续性得到。对固定正则支撑赋值，细分给出区间内容的有限可加性，穿孔正则性给出自上而下的右连续性。通常的 Lebesgue--Stieltjes 构造可直接扩张 远离目标 区间内容。在 趋向目标 一侧，单调边界残差
\[
 \alpha_s^\varphi=\lim_{b\downarrow0}U_\varphi(p_{T_s}(0,b))\ge0
\]
被赋给目标单点。由于到达目标的序列属于边界类而不是穿孔消失类，该原子可以保留下来。分离与测度扩张细节见附录~\ref{app:regular}。
\end{proof}

对固定的正则支撑赋值，$\lambda_m((a,b])$ 是模块 $m$ 中跨越区间 $(a,b]$ 所得到的价值。不同正则支撑赋值可以给同一通行赋予不同的有限权重，同时仍共同与原始次序相容。若密度或参数核存在，它们只是对某个特定相容支撑赋值施加的派生限制。

目标达成原子是在给定正则支撑赋值下，从侧 $s$ 到达预设目标时附加的边界价值。一次触及目标的完整经历还包括其穿孔趋向目标 与 远离目标 通行。例如，在支撑赋值 $\Lambda$ 下，从右侧距离 $b$ 连续跨越到左侧距离 $a$ 的通行价值为
\[
 \lambda_{T_R}^{\Lambda}([0,b])-\lambda_{A_L}^{\Lambda}((0,a])
 =\alpha_R^{\Lambda}+\lambda_{T_R}^{\Lambda,\circ}((0,b])-\lambda_{A_L}^{\Lambda}((0,a]).
\]
因此，一次物理跨越贡献两个有向通行分量以及一次右侧目标达成事件。

定义正则价格零空间
\[
 J_{\mathrm{reg}}
 =\left\{z\in E:
 \varphi\!\left(\rho(z+I)\right)=0
 \text{ 对每个 }\varphi\in\mathcal S_{\mathrm{reg}}
 \right\}.
\]

\begin{corollary}[一个共同标量化的点识别]\label{cor:common-scale}
在定理~\ref{thm:B} 下，以下三项等价：
\begin{enumerate}
 \item 归一化正则支撑族 $\mathcal S_{\mathrm{reg}}$ 为单点集；
 \item 归一化 Stieltjes 族 $\mathfrak M_{\mathrm{reg}}(\succeq)$ 为单点集；
 \item $J_{\mathrm{reg}}$ 在 $E$ 中满足 $\codim J_{\mathrm{reg}}=1$。
\end{enumerate}
当这些条件成立时，未归一化的正则相容支撑赋值形成一条正射线，因此在归一化之前，四个有向测度和达成原子可以在一个共同正比例因子意义下联合识别。
\end{corollary}

在固定有限物理网格上，可容许的空集收缩区间序列最终会变为空集，因此 PROC 中的穿孔连续性部分成为空洞条件。PVBC 却不一定冗余，因为有限维账户空间仍允许无界的路径重复次数和非阿基米德次序；注~\ref{rem:finitegrid-pvbc} 给出了明确例子。因此，在不受限路径域上的点识别仍由上述余维条件决定。对于有限经验设计，归一化可加补全的唯一性则是一个由实际有效观察限制决定的有限维多面体单点问题。

在定性概率中，单调连续性是从有限可加走向可数可加的经典路径 \citep{Villegas1964,ChateauneufJaffray1984}。本文的连续性类则由穿孔通行的消失生成。连续的正支撑赋值足以分离实际历史，而目标单点仍可作为一个边界原子保留。

\section{端点进展、总通行调整与目标达成}\label{sec:pha}

本节逐个解释每个正则支撑赋值，把由有序端点固定的部分，与附着于总遍历、对反向具有偶对称性的调整，以及目标达成产生的边界贡献分离开来。在定理~\ref{thm:B} 下，固定 $\Lambda\in\mathfrak M_{\mathrm{reg}}(\succeq)$，并在代数推导中省略其上标。该分解对每个赋值分别成立；下面的序数方案结果则针对原始次序或相容补全集合表述。

所有转折、返回和重复都通过第~\ref{sec:compression} 节定义的总剖面 $V_s^\gamma$ 进入；引理~\ref{lem:net} 则由有序端点识别相应的净剖面。

把穿孔限制记为
\[
 \lambda_{T_s}^{\circ}=\lambda_{T_s}|_{(0,R_s]},
 \qquad
 \lambda_{A_s}^{\circ}=\lambda_{A_s}|_{(0,R_s]},
\]
并定义
\[
 \mu_s=(\lambda_{T_s}^{\circ}+\lambda_{A_s}^{\circ})/2,
 \qquad
 \kappa_s=(\lambda_{T_s}^{\circ}-\lambda_{A_s}^{\circ})/2.
\]
于是 $\kappa_s\ll\mu_s$，且 Radon--Nikodym 导数
\[
 \theta_s=\dd\kappa_s/\dd\mu_s
\]
存在，并在 $\mu_s$ 几乎处处满足 $|\theta_s|\le1$。令 $P(0)=0$，并对侧 $s$ 上的 $x$ 定义
\[
 P(x)=-\mu_s((0,|x|]).
\]
再令 $A_s(\gamma)=h^{T_s}_\gamma(0)$ 表示从侧 $s$ 的目标达成次数。

\begin{theorem}[势--总通行--达成分解]\label{thm:C}
对每个归一化、正则且相容的支撑赋值 $\Lambda$，令 $P_\Lambda(0)=0$。其通行价值可以分解为
\[
\boxed{
 U_\Lambda(\gamma)
 =P_\Lambda(\gamma(1))-P_\Lambda(\gamma(0))
 +\sum_{s\in\{L,R\}}\int V_s^\gamma\dd\kappa_{s,\Lambda}
 +\sum_{s\in\{L,R\}}\alpha_{s,\Lambda} A_s(\gamma).}
\]
三项依次为端点进展、总通行调整和目标达成分量。等价地，中间项为
\[
 \sum_s\int V_s^\gamma\theta_{s,\Lambda}\dd\mu_{s,\Lambda}.
\]
对每个固定的 $\Lambda$，其归一化四个有向测度与 $(P_\Lambda,\kappa_{L,\Lambda},\kappa_{R,\Lambda},\alpha_{L,\Lambda},\alpha_{R,\Lambda})$ 之间的变换可逆。在推论~\ref{cor:common-scale} 下，归一化分解得到点识别；否则，本定理逐个赋值描述整个相容族。
\end{theorem}

\begin{proof}
在每一侧，
\[
\begin{aligned}
 \int N_s^T\dd\lambda_{T_s}^{\circ}
 -\int N_s^A\dd\lambda_{A_s}^{\circ}
 &=\int D_s^\gamma\dd\mu_s+\int V_s^\gamma\dd\kappa_s.
\end{aligned}
\]
引理~\ref{lem:net} 把第一项积分转化为 $P(\gamma(1))-P(\gamma(0))$。目标原子贡献 $\alpha_s A_s(\gamma)$。对固定支撑赋值，$\mu_s$、$\kappa_s$ 与 $\alpha_s$ 的定义均可逆，再结合 $P(0)=0$，就得到所述一一重参数化。
\end{proof}

这些名称的区分具有实质内容。有符号测度 $\kappa_s$ 是\emph{总通行调整测度}：它记录在剔除端点分量后仍然保留的、对反向具有偶对称性的共同溢价或惩罚。由于 $V_s^\gamma$ 在一条开放单调路径上通常并非零，中间泛函是一种一般的总通行调整。对于穿越区间 $E$ 的闭合目标外循环，端点贡献为零，剩余价值 $2\kappa_s(E)$ 称为\emph{循环通行价值}。这里“通行迟滞”一词仅指这种静态、反转偶对称的调整；动态迟滞是另一种结构。

\begin{corollary}[直接、绕行与触及目标历史]\label{cor:three-history}
固定 $0<a<b<c\le R_R$ 以及右侧历史
\[
 \gamma^D:c\to b,
 \qquad
 \gamma^C:c\to a\to b,
 \qquad
 \gamma^A:c\to0\to b.
\]
它们的完整分解为
\[
\begin{aligned}
 U(\gamma^D)
  &=P(b)-P(c)+\kappa_R((b,c]),\\
 U(\gamma^C)
  &=P(b)-P(c)+\kappa_R((b,c])+2\kappa_R((a,b]),\\
 U(\gamma^A)
  &=P(b)-P(c)+\kappa_R((b,c])+2\kappa_R((0,b])+\alpha_R.
\end{aligned}
\]
因此
\[
\begin{aligned}
 U(\gamma^C)-U(\gamma^D)
  &=2\kappa_R((a,b]),\\
 U(\gamma^A)-U(\gamma^D)
  &=2\kappa_R((0,b])+\alpha_R.
\end{aligned}
\]
第一个差分是额外目标外“减量--复发”绕行的循环通行价值；第二个差分是额外触及目标绕行的总通行调整，加上一次从右侧到达目标的价值。左侧具有完全类似的公式。
\end{corollary}

\begin{proof}
总通行剖面为
\[
\begin{aligned}
 V_R^{\gamma^D}&=\1_{(b,c]},\\
 V_R^{\gamma^C}&=\1_{(b,c]}+2\1_{(a,b]},\\
 V_R^{\gamma^A}&=\1_{(b,c]}+2\1_{(0,b]}.
\end{aligned}
\]
三条历史都有有序端点 $(c,b)$，故其端点贡献都是共同项 $P(b)-P(c)$。只有 $\gamma^A$ 到达目标，并且恰好从右侧到达一次。代入定理~\ref{thm:C} 即得三个完整公式及两个差分。
\end{proof}

这一重参数化在代数上很简单；其解释力来自路径几何对端点、重复通行和目标达成分量的分离。

\subsection{无需选定赋值的反事实方案转移}

上述分解为每个支撑赋值提供了基数描述，但可加通行记账还给出一个纯序数限制，而后者对受控方案选择更有用。只要两个实际历史对产生相同的账户差分，它们就必须具有相同的原始排名。该结论既不需要点识别，也不需要从支撑族中选定某一个赋值。

\begin{theorem}[实际比较对差分不变性与方案转移]\label{thm:D}
假设 BPCC。若四条实际历史满足
\[
 q(\gamma)-q(\eta)=q(\gamma')-q(\eta'),
\]
则
\[
 \gamma\succeq\eta
 \quad\Longleftrightarrow\quad
 \gamma'\succeq\eta'.
\]
严格偏好与无差异也满足同样的等价关系。

特别地，固定 $a>0$，并考虑两个右侧方案环境 $j=1,2$，满足 $0<a<b_j<c_j$。定义
\[
 \gamma_j^A:c_j\to0\to b_j,
 \qquad
 \gamma_j^C:c_j\to a\to b_j.
\]
则
\[
 q(\gamma_1^A)-q(\gamma_1^C)
 =q(\gamma_2^A)-q(\gamma_2^C),
\]
因此，在两个环境中，触及目标的方案相对于在 $a$ 处转折的方案具有完全相同的弱偏好、严格偏好或无差异关系。若每一对内部的非通行后果都得到控制，则即使两个环境的初始状态和终止状态不同，一个已观察的方案排名也能预测另一个。
\end{theorem}

\begin{proof}
由命题~\ref{prop:realcone}，在 BPCC 下，一个实际差分属于揭示锥 $C$ 当且仅当对应的弱比较成立。因此，相同账户差分意味着两个弱比较等价。对反向差分应用同样论证，又得到反向弱比较等价，从而严格偏好和无差异也等价。

对于所展示的目标接触方案，共同部分相互抵消。相对于 $\gamma_j^C$，路径 $\gamma_j^A$ 增加一次在 $[0,a]$ 上到达目标的趋向目标 通行，并把 $(a,b_j]$ 上的远离目标 通行替换成 $(0,b_j]$ 上的远离目标 通行。最终账户差分对两个 $j$ 都是 $(0,a]$ 上相同的“到达目标/远离目标”对比，与 $b_j$、$c_j$ 无关。
\end{proof}

在每一个正则支撑赋值下，定理~\ref{thm:D} 中的共同差分具有价值
\[
 2\kappa_{R,\Lambda}((0,a])+\alpha_{R,\Lambda}.
\]
这一恒等式在整个支撑族上解释了被转移的比较。仅依赖端点的通行偏好把该共同对比设为零，因此在两个菜单中都预测无差异。于是，一个菜单中的严格受控选择会拒绝该子类；在 BPCC 下，它还会预测每一个能够实现同一通行账户差分的其他实际菜单中的相应严格选择。

\begin{corollary}[目标外循环与收缩的目标返回循环]\label{cor:loops}
对侧 $s$ 上穿过 $(a,b]$ 的典范目标外往返，
\[
 U(C_{s;a,b})=2\kappa_s((a,b]).
\]
对一个半径为 $\varepsilon$、从目标出发进入侧 $s$ 并从该侧返回目标的循环，
\[
 U(C_{\varepsilon,s})
  =2\kappa_s((0,\varepsilon])+\alpha_s
  \longrightarrow\alpha_s
 \qquad\text{当 }\varepsilon\downarrow0.
\]
因此，在任意固定支撑赋值下，目标外循环隔离总通行调整测度；而收缩的目标返回循环在穿孔调整消失后隔离目标达成残差。有限序数数据通过整个相容族识别这些量，而不是仅凭某一个赋值层面的恒等式。
\end{corollary}

\begin{corollary}[仅依赖端点的子类]\label{cor:endpoint-class}
在定理~\ref{thm:B} 的完整路径域上，原始通行偏好仅依赖有序端点，当且仅当对每个 $\Lambda\in\mathfrak M_{\mathrm{reg}}(\succeq)$ 都有
\[
 \kappa_{L,\Lambda}=\kappa_{R,\Lambda}=0,
 \qquad
 \alpha_{L,\Lambda}=\alpha_{R,\Lambda}=0.
\]
左右两侧仍可以具有不同的势测度；除非推论~\ref{cor:common-scale} 成立，不同的正则支撑赋值也可以导出不同的以目标为中心的端点势。
\end{corollary}

\begin{proof}
若原始通行偏好仅依赖有序端点，则每一个目标外往返和每一个目标返回循环都与停留无差异。精确一致同意于是使每个支撑赋值在这些循环上取零。推论~\ref{cor:loops} 先给出所有穿孔区间上的 $\kappa_s=0$，继而得到 $\alpha_s=0$。反过来，若每个支撑赋值都满足 $\kappa_L=\kappa_R=0$ 且 $\alpha_L=\alpha_R=0$，则定理~\ref{thm:C} 使每个支撑价值仅依赖有序端点；定理~\ref{thm:B} 的一致同意进一步把同一性质传递给原始次序。
\end{proof}

在固定初始状态下，这一子类就是一个以目标为中心的终止状态次序。定理~\ref{thm:C} 隔离了由总有向通行和实际目标达成所携带的额外历史内容。

\section{匹配方案与有限数据预测}\label{sec:implications}

表示结果描述完整的条件通行次序，但经济应用通常只能观察有限多个方案选择。本节直接处理这些观察。主要结果从序数方案数据出发，为反事实通行价值给出锐界；随后用一个简短的货币校准推论说明，同一有限设计几何如何以货币单位表达。

\subsection{来自序数方案选择的锐界}\label{sec:finiteprediction}

固定一个有限物理网格 $\Pi$。令 $x_1,\ldots,x_J$ 表示该网格上的有利单单元通行差分：趋向目标 单元按“通行相对于停留”的方向取向，远离目标 单元按“停留相对于通行”的方向取向。它们张成有限维网格账户空间 $E_\Pi$。令
\[
 e_\Pi=\sum_{j=1}^J x_j.
\]
设 $\mathcal O=\{(\gamma_i,\eta_i)\}_{i=1}^n$ 为网格上已观察的弱方案选择，并假定每个观察对内部维持的背景属性得到控制；记 $d_i=q(\gamma_i)-q(\eta_i)$。定义
\[
 \mathcal P_\Pi(\mathcal O)
 =\left\{
 p\in E_\Pi^*:
 p(e_\Pi)=1,\quad
 p(x_j)\ge0\ \forall j,\quad
 p(d_i)\ge0\ \forall i
 \right\}.
\]
这些弱单元不等式给出有限设计的闭合、面向目标的补全集合。

对同一网格上的一个未观察方案对 $(\gamma,\eta)$，令 $d=q(\gamma)-q(\eta)$，并定义
\[
 \underline\Delta_{\mathcal O}(d)
 =\min_{p\in\mathcal P_\Pi(\mathcal O)}p(d),
 \qquad
 \overline\Delta_{\mathcal O}(d)
 =\max_{p\in\mathcal P_\Pi(\mathcal O)}p(d).
\]

\begin{theorem}[有限设计反事实通行锐界]\label{thm:finite-bounds}
假设 $\mathcal P_\Pi(\mathcal O)\neq\varnothing$。则它是一个紧凸多面体，并且对每个反事实差分 $d$，
\[
 \{p(d):p\in\mathcal P_\Pi(\mathcal O)\}
 =\bigl[\underline\Delta_{\mathcal O}(d),\overline\Delta_{\mathcal O}(d)\bigr].
\]
两个端点分别由两个线性规划问题求得。因此，严格为正的下界稳健地把 $\gamma$ 排在 $\eta$ 之上；严格为负的上界稳健地把 $\eta$ 排在 $\gamma$ 之上；若区间同时包含正负值，则记录了相容可加补全之间的反事实歧义。
\end{theorem}

\begin{proof}
因为 $p(x_j)\ge0$ 且 $p(e_\Pi)=\sum_jp(x_j)=1$，每个单元价值都位于 $[0,1]$。又由于 $x_j$ 张成 $E_\Pi$，归一化泛函有界；其余限制只是有限多个闭线性不等式。因此 $\mathcal P_\Pi(\mathcal O)$ 是紧凸多面体。在映射 $p\mapsto p(d)$ 下，其线性像为一个紧区间，端点由上述两个线性规划达到。多面体中的每个元素都定义有限观察的一个弱面向目标可加补全，由此得到锐性。
\end{proof}

定理~\ref{thm:D} 是精确转移的特例：当一个已观察比较对与一个反事实比较对具有相同的通行账户差分时，同一可加限制同时适用于二者。当不存在任何一个单独的已观察比较拥有目标差分时，有限数据定理还可以组合多个绕行、反向和目标接触选择的信息。

\subsection{货币校准}\label{sec:wtp}

有限的“方案--价格”数据可以在不为完整原始次序选定单一支撑赋值的情况下，把同样的可加对比校准为货币单位。令一个观察对象为 $(\gamma,m)$，并在该有限设计内维持货币拟线性。相容的可加通行赋值满足
\[
 p\bigl(q(\gamma_i)-q(\eta_i)\bigr)\ge m_i-n_i
\]
对每一个已观察弱选择 $(\gamma_i,m_i)\succeq(\eta_i,n_i)$ 成立，并同时满足所选网格上的目标方向不等式。

\begin{corollary}[货币校准的通行锐界]\label{cor:wtp}
对任意反事实通行差分 $d$，在相容有限设计多面体上，货币校准值集合 $\{p(d)\}$ 是一个区间，可能无界；其有限端点由两个线性规划求得。若相对于 $\eta$，方案 $\gamma$ 还需要支付额外费用 $c$，则严格高于 $c$ 的下界给出对 $\gamma$ 的稳健严格偏好，而严格低于 $c$ 的上界给出对 $\eta$ 的稳健严格偏好。
\end{corollary}

该推论只是对有限设计识别集合的货币校准。本文的经济比较并不依赖全局拟线性补全，方案--价格计算也没有被用于完整定义域的表示定理。

\subsection{识别与适用范围}

固定有限物理网格会消除非平凡的穿孔收缩序列，因此 PROC 正则性问题在该网格上成为空洞条件。当不受限路径域允许无界重复时，PVBC 仍有实质内容：注~\ref{rem:finitegrid-pvbc} 给出了一个满足有限消去与目标方向、但违反 PVBC 的字典序有限网格次序。相比之下，有限经验设计只使用有限多个已观察不等式，其可加补全问题按构造就是多面体问题。

在序数方案数据下，与决策相关的识别对象通常是
\[
 [\underline\Delta_{\mathcal O}(d),\overline\Delta_{\mathcal O}(d)].
\]
在方案--价格数据下，同一有限设计不等式把这个对比区间校准到货币单位。无论是序数区间还是货币校准区间，都可以在不分别识别 $\alpha_s$、$\kappa_s$ 或 $P$ 的情况下直接用于决策。精确无差异可以使序数相容多面体收缩。在一个 $d$ 维设计上，$d-1$ 个独立无差异等式加上归一化足以固定唯一可行的可加补全；更一般地，单点识别应根据全部有效约束系统检查。

条件通行次序最直接地由受控方案比较观察得到。当维持的背景属性被匹配或单独净除时，序数转移与归一化有限网格界均适用。方案--价格菜单提供一种明确方式来放松“货币成本相同”的要求：费用通过拟线性计价物进入，而通行分量仍保持集合识别。其他健康、持续时间和动态后果，在其价值已知或已建模时，可以通过应用特定结构加入。

BPCC 仍是不受限路径域上的一个全局有限列表限制。附录~\ref{app:primitive} 给出一个六历史族，其中原始平衡证书需要无界的重复次数。在有限经验设计中，可加可行性则直接通过相应的线性不等式系统检验。

\section{结论}\label{sec:conclusion}

本文构造了一种基于通行的一维目标相对历史压缩，并在所得账户空间上建立可加测量理论。通用载体识别可加对象，精确纤维定理则刻画保持该对象不变的可行时间顺序变换。在有限步一维定义域中，有序端点与总跨越剖面随后可以重构完整有向账本。

BPCC 把实际历史比较转化为精确的偏可加次序。PVBC 刻画何时阿基米德完备化保持实际历史次序，从而允许用归一化实数支撑赋值的一致同意作精确表示。有限网格的字典序例子说明：即使物理状态网格有限，只要路径可无界重复，这一限制仍具有实质内容。PROC 提供一个对边界敏感的正则性细化，从而得到允许目标原子的有向 Stieltjes 测度。

测度表示逐个赋值给出“端点--总通行--目标达成”的解释。经济学内容则在原始排名与识别集合层面表述。对于任意有限个具有共同端点、且在应用上可行的遍历，只要匹配时间安排足够长，就可以仅通过增加端点停留，在不改变遍历速度的情况下固定持续时间和一个能够区分端点的连续累计流量统计量。相同账户差分可以跨环境转移序数排名，有限个已观察选择则对反事实通行对比给出锐界；方案--价格数据还可以把同一界表示为货币单位。

通行账户有意把可交换的闭合游程视为在时间顺序上中性。纤维定理使这一限制精确化：它识别压缩究竟丢失了哪些信息，而不是把所有终点相同的历史都视为等价。若要评价闭合游程的顺序、首次达成效应、内生目标或动态状态转移，则需要更丰富的记录；这些也构成超越本文可加通行域的自然扩展。

\section{术语、记号与适用范围约定}\label{app:notation}

本附录是一份读者地图，而不是所有证明局部符号的清单。表~\ref{tab:dependency} 记录表示理论的逻辑主干；表~\ref{tab:notation} 汇集在多个正文小节中反复使用的术语。只在单个证明中出现的符号，则在其出现位置定义。

\begin{table}[ht]
\centering
\small
\renewcommand{\arraystretch}{1.16}
\begin{tabularx}{\textwidth}{@{}>{\raggedright\arraybackslash}p{0.39\textwidth}>{\raggedright\arraybackslash}X@{}}
\toprule
条件或继承结果 & 在表示链中的作用\\
\midrule
有限步路径域；弱时间变换不变性；拼接可加性 & 所有允许的群值通行赋值都必须通过的通用载体（定理~\ref{thm:universal}）。\\
一维带标签路径几何 & 载体的精确可行纤维（定理~\ref{thm:fiber}）。\\
有限步一维域中的有序端点 $+$ 总通行 & 完整有向账本的无损重构（命题~\ref{prop:gross-endpoints}）。\\
弱序 $+$ BPCC & 实际历史上的精确典范偏可加次序（命题~\ref{prop:envelope}）。\\
$+$ 面向目标的方向单调性 $+$ PVBC & 由归一化实数支撑赋值的一致同意作精确表示，同时无需将形式差分全序化（定理~\ref{thm:A}）。\\
$+$ 区间丰富性 $+$ PROC & 正则支撑赋值足以表示实际次序；每个赋值诱导四个有向 Stieltjes 测度以及各侧目标原子（定理~\ref{thm:B}）。\\
$+$ $\codim J_{\mathrm{reg}}=1$ & 一个归一化共同标量化的点识别（推论~\ref{cor:common-scale}）。\\
定理~\ref{thm:B} $+$ 净通行守恒 & 每个支撑赋值的端点--总通行--达成分解（定理~\ref{thm:C}）。\\
共同端点 $+$ 固定可行遍历 $+$ 允许在初始/终止端点停留 $+$ 一个区分端点的连续流量统计量 & 在足够长的共同时间跨度上，不改变遍历速度即可匹配持续时间与流量（命题~\ref{prop:matched-design}）。\\
BPCC $+$ 相同的实际通行账户差分 & 无需选定赋值即可转移方案的弱/严格/无差异序数排名（定理~\ref{thm:D}）。\\
有限经验网格 $+$ 已观察弱方案选择 & 反事实价值锐界与稳健的样本外方案预测（定理~\ref{thm:finite-bounds}）。\\
有限方案--价格设计 $+$ 拟线性计价物 & 通行对比的货币校准锐界（推论~\ref{cor:wtp}）。\\
\bottomrule
\end{tabularx}
\caption{主要表示结果及其核心含义之间的依赖关系图。}
\label{tab:dependency}
\end{table}

\begingroup
\small
\renewcommand{\arraystretch}{1.16}
\setlength{\LTpre}{0.4em}
\setlength{\LTpost}{0.6em}
\begin{longtable}{@{}>{\raggedright\arraybackslash}p{0.20\textwidth}>{\raggedright\arraybackslash}p{0.25\textwidth}>{\raggedright\arraybackslash}p{0.46\textwidth}@{}}
\caption{核心术语、记号与适用范围约定}\label{tab:notation}\\
\toprule
类别 & 术语或符号 & 含义及维持的约定\\
\midrule
\endfirsthead
\multicolumn{3}{l}{\small\itshape 表~\thetable（续）}\\
\toprule
类别 & 术语或符号 & 含义及维持的约定\\
\midrule
\endhead
\midrule
\multicolumn{3}{r}{\small\itshape 下页续}\\
\endfoot
\bottomrule
\endlastfoot

\multicolumn{3}{@{}l}{\textit{经济对象与路径域}}\\*
条件通行偏好 & $\succeq$；$\{U_\Lambda\}_{\Lambda\in\mathfrak M_{\mathrm{reg}}}$ & 在非通行后果保持固定的条件下，个体对目标相对通行账户的主观排序；该排序由正则支撑赋值的一致同意恢复。\\
支撑赋值 & $U_\Lambda$ & 一个与原始次序相容的归一化可加泛函。除非正则族为单点集，单个 $U_\Lambda$ 不必保持每个原始严格比较，因此不能单独视为效用表示。\\
受控方案菜单 & --- & 被比较的一对方案中非通行后果相同的菜单，因此观察到的完整选择直接揭示条件通行次序。\\
比较路径域 & $\Pth$ & 用于表示的丰富有限步定义域；具体经验方案菜单可以只是其一个小子集。\\
目标相对状态 & $x_t$，目标 $0$ & 相对于预先给定目标测量的一维状态；在吸烟应用中为 $x_t=A_t-A_t^*$ 或保持次序的归一化。\\
侧与方向 & $s\in\{L,R\}$；$T,A$ & $L/R$ 表示目标以下/以上两侧；趋向目标 减少目标距离，远离目标 增加目标距离。这些都是几何标签。\\
目标达成 & $A_s(\gamma)$ & 从侧 $s$ 通过进入目标的趋向目标 通行到达 $0$。维持的原子按每次达成计值。\\
目标界面 & --- & 进入目标的趋向目标 单元包含 $0$，离开目标的远离目标 单元排除 $0$；一次连续跨越产生两个有向片段和一次达成。\\
有限步路径 & $\gamma:[0,1]\to X$ & 在有限多个单元上单调的连续路径；允许常值单元。\\
弱时间变换 & --- & $[0,1]$ 上连续非减满射；允许改变速度和加入停留，但保持访问状态的顺序。\\

\multicolumn{3}{@{}l}{\textit{通行账户与典范次序}}\\*
抽象通行账户 & $q(\gamma)\in\mathfrak L$ & 由模块--区间生成元在细分关系下取商所得的通用可加载体。\\
具体账本 & $\ell(\gamma)=(h^m_\gamma)_{m\in\M}$ & 记录所处一侧、方向、跨越水平、重复次数和目标达成的阶梯函数表示。定理~\ref{thm:universal} 将其与 $q(\gamma)$ 识别。\\
通行等价 & $\gamma\approx_p\eta$ & 可通过定理~\ref{thm:fiber} 中保持账户的操作相连的历史；等价地，$q(\gamma)=q(\eta)$。\\
BPCC & 平衡通行循环一致性 & 排除有限个弱改进（至少一个严格）组成、且总通行账户恰好平衡的列表。\\
典范包络 & $G_{\succeq}$，$J_{\succeq}$ & 由揭示账户差分生成的偏序商群及其典范路径映射。\\
PVBC & 成对消失基准一致性 & 要求阿基米德完备化保持实际历史比较；它排除实际的无穷小严格差分，同时不把形式差分全序化。\\
归一化支撑赋值 & $\mathcal S$ & 阿基米德化序单位空间上的正线性泛函，并在一个全支撑基准上归一化。在定理~\ref{thm:A} 下，它们的一致同意表示实际次序。\\

\multicolumn{3}{@{}l}{\textit{连续性与有向测度}}\\*
穿孔拓扑 & $\tau^\circ$ & 使收缩到空集的典范非目标通行收敛到零的最细局部凸拓扑；到达目标的序列被排除在外。\\
PROC & 穿孔揭示次序闭性 & 要求揭示锥的穿孔拓扑闭包不改变实际历史比较。\\
正则支撑赋值 & $\mathcal S_{\mathrm{reg}}$ & $\tau^\circ$ 连续的归一化支撑赋值；等价地，其价值在每个可容许穿孔收缩序列上趋于零。\\
有向测度 & $\lambda_{T_s}^{\Lambda},\lambda_{A_s}^{\Lambda}$ & 正则支撑赋值 $\Lambda$ 诱导的有限非负 Borel 测度；远离目标 的目标质量归一化为零。\\
穿孔测度 & $\lambda_{T_s}^{\Lambda,\circ},\lambda_{A_s}^{\Lambda,\circ}$ & 正则支撑赋值的有向测度在正目标距离 $(0,R_s]$ 上的限制。\\
达成系数与原子 & $\alpha_s^\Lambda$；$\alpha_s^\Lambda\delta_0$ & 支撑赋值特定的边界残差及其单点测度形式；可以为零，也不必在不同赋值之间相同。\\
\multicolumn{3}{@{}l}{\textit{有限数据方案选择}}\\*
方案--价格补全 & $v_p(\gamma,m)=p(q(\gamma))-m$ & 用于在有限网格上把通行价值校准到货币单位的应用特定拟线性扩展。\\
支付意愿区间 & $[\underline w(d),\overline w(d)]$ & 在与已观察方案--价格选择一致、面向目标的可加补全中，一个反事实方案差分所对应货币价值的锐集合。\\

\multicolumn{3}{@{}l}{\textit{分解与识别}}\\*
净通行 & $D_s^\gamma=N_s^T-N_s^A$ & 由有序端点决定的方向不平衡。\\
总通行 & $V_s^\gamma=N_s^T+N_s^A$ & 趋向目标 加 远离目标 的重复次数，进入反转偶对称调整与信息结果。\\
势测度与势 & $\mu_s$；$P$ & $\mu_s=(\lambda_{T_s}^{\circ}+\lambda_{A_s}^{\circ})/2$ 生成端点进展，且 $P(0)=0$。\\
总通行调整 & $\kappa_s$；$\int V_s^\gamma\,\dd\kappa_s$ & 有符号、反转偶对称分量，其中 $\kappa_s=(\lambda_{T_s}^{\circ}-\lambda_{A_s}^{\circ})/2$。\\
循环通行价值 & $2\kappa_s(E)$ & 闭合目标外往返通过 $E$ 时，在端点项抵消后的总通行调整。\\
目标返回循环 & $C_{\varepsilon,s}$ & 从目标出发到半径 $\varepsilon$ 再从侧 $s$ 返回；其价值收敛到 $\alpha_s$。\\
仅端点子类 & --- & 等价于 $\kappa_L=\kappa_R=0$ 且 $\alpha_L=\alpha_R=0$；左右势测度仍可以不同。\\

\end{longtable}
\endgroup

\input{zh/part12.tex}
\input{zh/part13.tex}
\input{zh/part14.tex}
\input{zh/part15.tex}
\input{zh/part16.tex}
\input{zh/part17.tex}

\end{document}
